% Figure: plan view of a pair of mitre lock gates closed against an upstream head,
% with the leaf thrust resolved into the mitre thrust along the leaf and the quoin
% (hinge-post) reaction. The two leaves meet at the mitre point M and lean upstream
% so the water load wedges them shut. No em-dashes.
\begin{figure}[htbp]
\centering
\begin{tikzpicture}[scale=1.0]
% chamber walls (plan view): two vertical lines bounding the lock chamber
\draw[enggray, very thick] (-3.6,-2.6) -- (-3.6,2.6);
\draw[enggray, very thick] ( 3.6,-2.6) -- ( 3.6,2.6);
\node[enggray, font=\scriptsize, anchor=east] at (-3.6,2.4) {chamber wall};
% quoin posts at the wall (hinge posts), one on each side
\fill[engblack] (-3.6,0) circle (2.2pt);
\fill[engblack] ( 3.6,0) circle (2.2pt);
\node[engblack, font=\scriptsize, anchor=south east] at (-3.6,0.05) {quoin $Q_1$};
\node[engblack, font=\scriptsize, anchor=south west] at ( 3.6,0.05) {quoin $Q_2$};
% high water (upstream) is below the gates in this plan; shade it
\fill[engblue!12] (-3.6,-2.6) rectangle (3.6,-0.05);
\node[engblue, font=\scriptsize] at (0,-2.1) {upstream (high) water};
\node[engblue, font=\scriptsize] at (0,1.9) {downstream (low) water};
% two leaves meeting at the mitre point M, pointing toward the high water (downstream apex)
% left leaf: from quoin (-3.6,0) to mitre point (0,-0.9)
% right leaf: from quoin (3.6,0) to mitre point (0,-0.9)
\coordinate (M) at (0,-0.9);
\draw[engblack, ultra thick] (-3.6,0) -- (M);
\draw[engblack, ultra thick] ( 3.6,0) -- (M);
\fill[engred] (M) circle (2.0pt);
\node[engred, font=\scriptsize, anchor=south] at (0,-0.82) {mitre $M$};
% water thrust on each leaf: normal to the leaf, pushing toward downstream (up)
\draw[engamber, ultra thick, ->] (-1.8,-0.45) -- ($(-1.8,-0.45)+(0.22,0.88)$);
\node[engamber, font=\scriptsize, anchor=east] at (-1.95,-0.2) {$P$};
\draw[engamber, ultra thick, ->] ( 1.8,-0.45) -- ($( 1.8,-0.45)+(-0.22,0.88)$);
\node[engamber, font=\scriptsize, anchor=west] at ( 1.95,-0.2) {$P$};
% mitre thrust (along leaf, compressive at M) and quoin reaction (along leaf, at Q)
\draw[englime, thick, ->] (M) -- ($(M)+(-0.95,0.24)$);
\node[englime, font=\scriptsize, anchor=north] at ($(M)+(-0.95,0.18)$) {$T_m$};
\draw[englime, thick, ->] (M) -- ($(M)+(0.95,0.24)$);
\node[englime, font=\scriptsize, anchor=north] at ($(M)+(0.95,0.18)$) {$T_m$};
\draw[engblue, thick, ->] (-3.6,0) -- ($(-3.6,0)+(0.95,-0.24)$);
\node[engblue, font=\scriptsize, anchor=south] at ($(-3.6,0)+(0.7,-0.18)$) {$R_q$};
\draw[engblue, thick, ->] ( 3.6,0) -- ($( 3.6,0)+(-0.95,-0.24)$);
\node[engblue, font=\scriptsize, anchor=south] at ($( 3.6,0)+(-0.7,-0.18)$) {$R_q$};
% mitre half-angle beta marked at M on the left leaf
\draw[engblack, thin] (M) -- ($(M)+(-1.3,0)$);
\node[engblack, font=\scriptsize, anchor=north east] at ($(M)+(-0.7,-0.02)$) {$\beta$};
\end{tikzpicture}
\caption[Plan of a mitre lock gate]{Plan view of a closed pair of mitre
gates. Each leaf runs from its quoin post at the chamber wall ($Q_1$, $Q_2$)
to the mitre point $M$, with the pair pointing into the high water so the
upstream head wedges them shut. The water thrust $P$ on a leaf (amber) acts
normal to the leaf face. Force balance on one leaf resolves $P$ into a
compressive mitre thrust $T_m$ (green) passing through the contact at $M$ and
a quoin reaction $R_q$ (blue) at the hinge post. The mitre half-angle
$\beta$ between a leaf and the line across the chamber sets the thrust
multiplication: the smaller $\beta$, the larger $T_m$ for a given $P$, which
is why lock leaves are set to a shallow mitre rather than square across.}
\label{fig:vol03ch10:lock-gate}
\end{figure}
